\section{Methodology}
\label{sec:methodology}

% =============================================================================
% SECTION SKELETON ONLY -- no prose yet. Each subsection below lists what it
% must contain. Agreed structure: 4.1 splits/sets -> 4.2 label regimes ->
% 4.3 model/training -> 4.4 metrics -> 4.5 experimental design -> 4.6 probing.
% Register: match the formal style of the Dataset section (notation table,
% explicit definitions). Interpretation (saturation, false-confidence, "which
% ruler won") is DEFERRED to Results/Discussion -- this section is design only.
% =============================================================================

Our experimental methodology builds on the label-bias and biased-ruler framework
that Parikh et al.~\cite{parikh2025labelbias,parikh2025whofails} developed for
breast DCE-MRI, carried over here to cervical-spine segmentation. We adopt
several of their design choices -- balancing the cohort by downsampling,
comparing expert (gold) against machine-generated (silver) reference labels, and
quantifying performance gaps with binarized rate-based fairness metrics -- and we
note throughout where the structure of CSpineSeg forces us to adapt them, chiefly
that each exam carries a gold or a silver label but never both.

\subsection{Data Splits and Label Sets}
\label{sec:splits}

The Dataset section left two pieces of notation undefined, marked $(\ast)$ in
\autoref{tab:notation}: the label sets and the train/validation/test split. We
close them here, and in doing so account for the drop from the working set
$\mathcal{D}_0$ ($N_0 = 1{,}254$ exams, whose composition is summarised in
\autoref{tab:cohort}) to the analysis cohort $\mathcal{D}$ ($N = 1{,}142$), the
two cohorts introduced in \autoref{tab:notation}.

Every exam carries exactly one reference segmentation, so the provenance flag
$q_i$ partitions the cohort into two disjoint sets,
\begin{equation}
  \mathcal{D}^{\mathsf{g}} = \{\, i : q_i = \mathsf{g} \,\},
  \qquad
  \mathcal{D}^{\mathsf{s}} = \{\, i : q_i = \mathsf{s} \,\},
  \qquad
  \mathcal{D}^{\mathsf{g}} \cap \mathcal{D}^{\mathsf{s}} = \varnothing,
\end{equation}
the expert (gold) set and the machine-generated (silver) set, with
$\mathcal{D} = \mathcal{D}^{\mathsf{g}} \cup \mathcal{D}^{\mathsf{s}}$. That an
image carries gold or silver but never both is a structural property of
CSpineSeg, and it shapes the ``\nameref{sec:design}''; we
return to it in ``\nameref{sec:regimes}''.

We partition the cohort at the patient level into training, validation, and test
sets $\mathcal{D}_{\mathrm{train}}, \mathcal{D}_{\mathrm{val}},
\mathcal{D}_{\mathrm{test}}$ in a $70/10/20$ ratio. Splitting by patient rather
than by exam keeps all of a patient's exams on the same side of every boundary,
so the handful of multi-exam patients cannot leak information from training into
evaluation. The split is stratified by the protected attributes $\mathcal{A}$
so that each subgroup's share stays roughly constant across
the three sets.\footnote{White and Black patients, who form the race comparison,
are stratified on the full race $\times$ age $\times$ sex key; the small
remaining race groups (\autoref{tab:cohort}) are too sparse to stratify and are
assigned to splits proportionally at random.}

On top of this we balance the sexes. The working set leans female
(\autoref{tab:cohort}), and an unequal base rate would confound the sex audit by
conflating a true performance gap with mere differences in representation, so
within each split we randomly downsample female exams -- the over-represented sex
throughout this cohort -- until the two are equal, removing $112$ in total
($N_0 = 1{,}254 \to N = 1{,}142$) and leaving every split exactly $50/50$
male/female. This removes representation as an explanation for any residual gap,
rather than attempting to close it~\cite{parikh2025labelbias}. The downsampling is seeded
and reproducible, and because it operates at the exam level while the
patient-level split already fixes which side each patient falls on, it introduces
no leakage even for the multi-exam patients. We adopt this balance-by-downsampling
protocol from Parikh et al.~\cite{parikh2025labelbias,parikh2025whofails}, who
balance their breast-MRI cohort by age to probe representational imbalance; we
apply the same idea to sex and to this dataset's multi-exam structure.
Stratification held in every split: race and age proportions match the cohort
marginals of \autoref{tab:cohort} to within a couple of percentage points, and
sex is exactly $50/50$ by construction. The
per-regime breakdown of these splits into gold and silver training cases is given
in ``\nameref{sec:regimes}''.

We keep the conventional $70/10/20$ split rather than a bare train/test partition, which
keeps the cohort partition portable across modelling choices. nnU-Net is
cross-validation-native -- it selects its configuration and post-processing through an
internal five-fold cross-validation rather than a held-out tuning
set~\cite{isensee2021nnunet} -- so for our segmentation models $\mathcal{D}_{\mathrm{train}}$
and $\mathcal{D}_{\mathrm{val}}$ are pooled into the single development set that the
cross-validation runs over, leaving none of the scarce expert annotations idle (most
valuable in the gold-only regime, ``\nameref{sec:regimes}''). The validation tier still
earns its place: a model that is not cross-validation-native -- a transformer baseline, or
any architecture used to extend this audit -- needs a held-out set for early stopping and
hyperparameter selection, so defining $\mathcal{D}_{\mathrm{val}}$ now lets such work tune
and report on the identical partition while $\mathcal{D}_{\mathrm{test}}$ remains the single
fully held-out test set for every method.

\subsection{Three Label Regimes}
\label{sec:regimes}

We want to know whether the provenance of the training labels -- expert
or machine-generated -- changes the fairness of the resulting model. To isolate
that effect we train three models that differ only in the labels they are given,
holding everything else fixed: the same architecture and training recipe
(``\nameref{sec:model}''), the same patient-level split of
``\nameref{sec:splits}'', and the same sex balance and stratification. Any
difference in their fairness behaviour is then attributable to the training
labels alone, and not to model capacity or cohort composition -- the assumption
that justifies every model-to-model comparison in this paper.

The three regimes follow directly from the disjoint gold/silver partition
$\mathcal{D} = \mathcal{D}^{\mathsf{g}} \cup \mathcal{D}^{\mathsf{s}}$ introduced
in ``\nameref{sec:splits}''. The mixed model $M_{\mathrm{mix}}$ trains on the
full training split $\mathcal{D}_{\mathrm{train}}$, taking each exam with
whichever label it carries; combining a small pool of expert labels with a
larger pool of automated ones is a standard response to the cost of expert
annotation in medical image segmentation~\cite{tajbakhsh2020embracing}, so
$M_{\mathrm{mix}}$ is the deployment-realistic model that a typical pipeline
would train. The
gold model $M_{\mathrm{gold}}$ and the silver model $M_{\mathrm{silver}}$
restrict training to only the expert or only the machine-generated exams, drawn
from $\mathcal{D}_{\mathrm{train}} \cap \mathcal{D}^{\mathsf{g}}$ and
$\mathcal{D}_{\mathrm{train}} \cap \mathcal{D}^{\mathsf{s}}$ respectively.
\autoref{fig:regimes} shows how the three regimes are drawn from the gold/silver
partition, and \autoref{tab:regimes} gives their full composition.

That the gold-only and silver-only regimes can be built at all rests on the
structural property of CSpineSeg flagged above: each exam carries a gold or a
silver label, never both ($\mathcal{D}^{\mathsf{g}} \cap \mathcal{D}^{\mathsf{s}}
= \varnothing$). This single-label structure is the chief way our setting departs
from the breast-MRI biased-ruler work we build on, and it shapes the experimental
design (``\nameref{sec:design}'').

\begin{figure}[t]
\centering
\begin{tikzpicture}
  \node[regnode=regimeCohort, text width=64mm] (D) at (0,2.5)
    {Analysis cohort $\mathcal{D}$\\ $1{,}142$ exams ($50/50$ M/F)\\[2pt]
     {\scriptsize each exam carries gold \emph{or} silver, never both}};

  \node[regnode=regimeGold,   text width=30mm] (DG) at (-2.7,0.7)
    {Gold set $\mathcal{D}^{\mathsf{g}}$\\[1pt] {\scriptsize expert annotation}};
  \node[regnode=regimeSilver, text width=32mm] (DS) at ( 2.7,0.7)
    {Silver set $\mathcal{D}^{\mathsf{s}}$\\[1pt] {\scriptsize machine-generated}};

  \node[regnode=regimeGold,   text width=22mm] (MG) at (-4.6,-1.8)
    {$M_{\mathrm{gold}}$\\[1pt] train $288$\\ {\scriptsize gold only}\\
     {\scriptsize ($50/50$ M/F)}};
  \node[regnode=regimeMix,    text width=38mm] (MM) at ( 0,  -1.8)
    {$M_{\mathrm{mix}}$\\ train $798$\\ {\scriptsize $318$ gold $+$ $480$ silver}\\
     {\scriptsize ($50/50$ M/F)}};
  \node[regnode=regimeSilver, text width=22mm] (MS) at ( 4.6,-1.8)
    {$M_{\mathrm{silver}}$\\[1pt] train $450$\\ {\scriptsize silver only}\\
     {\scriptsize ($50/50$ M/F)}};

  \draw[regedge=regimeCohort] (D)  -- (DG);
  \draw[regedge=regimeCohort] (D)  -- (DS);
  \draw[regedge=regimeGold]   (DG) -- (MG);
  \draw[regedge=regimeGold]   (DG) -- (MM);
  \draw[regedge=regimeSilver] (DS) -- (MS);
  \draw[regedge=regimeSilver] (DS) -- (MM);
\end{tikzpicture}
\caption{The three label regimes. The provenance flag $q_i$ partitions the
analysis cohort $\mathcal{D}$ into a disjoint gold (expert) set
$\mathcal{D}^{\mathsf{g}}$ and silver (machine-generated) set
$\mathcal{D}^{\mathsf{s}}$; the mixed model $M_{\mathrm{mix}}$ trains on both,
while $M_{\mathrm{gold}}$ and $M_{\mathrm{silver}}$ each train on one. Counts are
training exams and do not add up across regimes ($798 \neq 288 + 450$) because
each regime is sex-balanced independently; the full per-split composition is in
\autoref{tab:regimes}.}
\label{fig:regimes}
\end{figure}

\begin{table}[t]
\centering
\caption{The three label regimes. All three share the architecture, training
recipe, patient-level split, stratification, and exact $50/50$ sex balance, and
differ only in the provenance of their training labels. Counts are exams; every
split in every regime is exactly $50/50$ male/female. Validation and test counts
are regime-specific and are not comparable across rows.}
\label{tab:regimes}
\small
\begin{tabular*}{\textwidth}{@{\extracolsep{\fill}} l l r r r}
\hline\noalign{\smallskip}
Model & Training labels & Train & Val & Test \\
\noalign{\smallskip}\hline\noalign{\smallskip}
$M_{\mathrm{mix}}$    & gold $+$ silver & 798 (318 gold $+$ 480 silver) & 116 & 228 \\
$M_{\mathrm{gold}}$   & gold only       & 288                            &  44 &  76 \\
$M_{\mathrm{silver}}$ & silver only     & 450                            &  66 & 138 \\
\noalign{\smallskip}\hline
\end{tabular*}
\end{table}

Because the gold and silver sets are disjoint, each regime is sex-balanced on
its own, and the counts do not add up across regimes. The principle is that a
set can be exactly $50/50$ male/female while neither of its two parts is. The
mixed training split is balanced across all $798$ of its exams ($399$ female and
$399$ male), but that balance is a cancellation: its $318$ gold exams lean female
($174$ versus $144$) while its $480$ silver exams lean male ($255$ versus $225$),
so the $30$-exam female surplus on the gold side exactly offsets the $30$-exam
male surplus on the silver side. Re-balancing sex within the gold pool alone
therefore drops those $30$ surplus female exams, which is why $M_{\mathrm{gold}}$
trains on $288$ gold exams ($144$ of each sex) rather than the $318$ that appear
in the mixed split -- a balanced subset of the very same cases, not a different
draw. The silver regime mirrors it, dropping the $30$ surplus male exams to reach
$450$ ($225$ of each sex). The regimes consequently do not sum --
$288 + 450 \neq 798$ -- because each discards its own surplus sex independently.
Validation and test counts are likewise regime-specific (\autoref{tab:regimes})
and should not be compared across rows.

One asymmetry remains by construction: the gold training set ($288$ exams) is
smaller than the silver one ($450$), simply because fewer CSpineSeg exams were
expertly annotated in the first place. This reflects the dataset's structure
rather than a design choice, and we flag it as a limitation of the
gold-versus-silver comparison.

A second concern is subtler. Because $M_{\mathrm{gold}}$ and $M_{\mathrm{silver}}$
train on disjoint exams, a demographic imbalance between the gold and silver sets
would reintroduce cohort composition as an explanation for any fairness difference
between them. It does not arise here: Zhou et al.~chose which exams to
expert-annotate pseudo-randomly -- by medical-record number rather than by any
demographic criterion~\cite{zhou2025cspineseg} -- so the gold and silver sets are
demographically comparable, and label provenance is the only variable that
separates them.

\subsection{Model and Training}
\label{sec:model}

We segment every regime with nnU-Net~\cite{isensee2021nnunet}, the self-configuring
framework that has become the standard baseline for biomedical segmentation. Rather
than fix an architecture by hand, nnU-Net derives the whole pipeline -- patch and
batch size, network depth, resampling, and intensity normalisation -- from a
fingerprint of the dataset (image sizes, voxel spacings, modality, class
frequencies) through a set of rule-based heuristics. Two properties make it the right
backbone for a fairness audit. First, it is a genuinely strong baseline: under
controlled, equal-resource validation a well-scaled nnU-Net still matches or exceeds
the transformer and state-space architectures that claim to surpass
it~\cite{isensee2024revisited}, so any performance gap we report is a property of the
data and its labels, not an artefact of a weak model. Second, because the
configuration is derived from the fingerprint rather than chosen per run,
holding the framework fixed across the three regimes guarantees that
$M_{\mathrm{mix}}$, $M_{\mathrm{gold}}$, and $M_{\mathrm{silver}}$ share an identical
architecture and training recipe -- the assumption that justifies every model-to-model
comparison in ``\nameref{sec:regimes}''.

We use nnU-Net v2 with the residual-encoder ``L'' preset (ResEnc-L), the framework's
current recommended default, which devotes a larger GPU-memory budget to a
residual-encoder U-Net~\cite{isensee2024revisited}. Because CSpineSeg is MRI, the planner
selects per-case $z$-score intensity normalisation. Following nnU-Net's standard protocol
we train two configurations for every regime -- a 2D model and a full-resolution 3D model
(patches $512\times512$ and $15\times512\times512$, the latter's shallow depth matching
the thin, few-slice sagittal stacks of CSpineSeg) -- each as a five-fold
cross-validation over the pooled development set
($\mathcal{D}_{\mathrm{train}}\cup\mathcal{D}_{\mathrm{val}}$) defined in
``\nameref{sec:splits}''. nnU-Net's own
configuration search then selects, for each regime, the softmax-averaged ensemble of
both configurations across all five folds -- ten trained models -- as the final
predictor. The same recipe runs for all three regimes, changing only the training
labels.

Two choices depart from nnU-Net's defaults, both to serve the audit. We replace the
framework's random cross-validation folds with demographically stratified ones, so that
each fold preserves the joint race~$\times$~age~$\times$~sex composition of the cohort;
together with the patient-level, sex-balanced split of ``\nameref{sec:splits}'' this
keeps subgroup representation constant through training and validation. And we apply no
per-structure largest-component post-processing: nnU-Net's default ``keep largest
component'' heuristic is built for single-object targets and would collapse the
multi-vertebra, multi-disc anatomy of a cervical spine into a single blob -- on our
validation folds it cut vertebral-body Dice from $0.96$ to $0.49$ -- so we keep the raw
network output. Training otherwise uses the stock nnU-Net trainer; our only addition is
experiment logging, leaving the loss, optimiser, and augmentation untouched. The splits
and the sex-balancing downsample are seeded and reproducible.
% TODO: models not yet released -- add public-model link here once the Hugging Face
% repo is live (see docs/nnunet/07_huggingface_release.md). E.g.: ", and the trained
% models are released publicly at \url{...}." Do NOT claim release until it is public.
Training all thirty
model-folds consumed roughly $1{,}105$ GPU-hours
on A100 and L40S hardware; under the standard ML~CO$_2$
accounting~\cite{lacoste2019quantifying,strubell2019energy,patterson2021carbon} and the
2023 Danish grid intensity~\cite{dea_keyfigures2023} this corresponds to an estimated
$53$--$74$\,kg~CO$_2$e.

\subsection{Evaluation Metrics}
\label{sec:metrics}

% - Overlap/distance metrics: Dice and HD95 (95th-percentile Hausdorff, mm).
% - nDSC: define as the class-imbalance-normalised Dice of Raina et al. (2023)
%   -- normalises the false-positive penalty by a reference class/lesion load
%   so the score stops rewarding large structures. EXPLICIT tie-back: this is
%   the "volume-aware score introduced in the Methodology" that the Dataset
%   section (sec:confounders) promised as the volume-confound fix; name it.
% - Fairness metrics, binarized rate-based (canonical Parikh / Fairlearn /
%   EEOC): per-case beneficial outcome (success), per-group success rate
%   P(yhat=1 | group); DPD = max(rate) - min(rate); DIR = min/max.
%   Thresholds: Dice/nDSC 0.8, HD95 5 mm (configurable). Four-fifths rule
%   (DIR < 0.8 flags adverse impact). Sensitivity sweep emitted per run.
%   Portability caveat ("four-fifths from US employment law; adopted as field
%   convention + for comparability, not a legal claim").
% - Significance backbone: Mann-Whitney U (sex, race; effect size r_rb) and
%   Kruskal-Wallis H (age 3-bin; eps^2) on the CONTINUOUS scores; BH-FDR across
%   the family.
% - A-PRIORI primacy argument (do NOT justify via saturation -- that leaks a
%   result): at macro Dice ~0.89, well above the 0.8 beneficial-outcome cutoff,
%   the binarized DIR is near-degenerate by construction (few failures per
%   group -> low power), so the continuous tests are the more sensitive
%   instrument BY DESIGN, and the sweep guards the DIR reading. Frame as: we
%   report DIR for field-comparability + four-fifths, but the high-Dice regime
%   makes the continuous tests primary. (This also pre-motivates the sweep.)

\subsection{Experimental Design}
\label{sec:design}

% - Generated-silver ruler MECHANISM: train M_gold (D002), predict on the 76
%   gold-test images -> those predictions ARE the "silver ruler" (Ruler B);
%   expert labels are the "gold ruler" (Ruler A). Mirrors how the original
%   silver labels were made (gold-trained nnU-Net applied to unseen cases).
% - WHY generate it / why not reuse Zhou et al.'s silver model: (a) weights
%   unavailable; (b) their split may overlap our gold test cases (leakage). Our
%   M_gold is trained on our split, which excludes the test cases by design.
% - Note: scoring M_mix against its OWN predictions is degenerate (Dice = 1);
%   rejected. M_mix and M_gold are independently trained, so Dice != 1.
% - E1-E4 design TABLE, stated NEUTRALLY (no saturation / false-confidence /
%   "which ruler won" -- those are Results/Discussion):
%       E1 Global audit       M_mix          | full 228 test, mixed labels
%       E2 Biased ruler       M_mix          | 76 gold test, gold vs gen-silver
%       E3 Bias amplification M_gold vs M_silver | 76 gold test, gold labels
%       E4 Mixed vs gold      M_mix vs M_gold     | 76 gold test, gold labels
% - MAKE THE TEST-SET DISTINCTION EXPLICIT: E1 is on the full 228 against mixed
%   labels; E2-E4 all sit on the SAME 76-case gold test set with expert labels
%   (apples-to-apples; differences attributable to labels/training, not data).
% - E4 is the production-realistic companion to E3 (same axis: does silver in
%   training degrade fairness?). Keep it as a row here for design transparency;
%   Results folds it into the E3 narrative rather than a standalone finding.

\subsection{Encoder Probing}
\label{sec:probing}

% - Short, last on purpose: diagnostic, not part of the core train/eval loop.
% - Setup: extract features from M_mix's bottleneck vs a generic medical-image
%   encoder (MRI-CORE); PCA + UMAP for visualisation; linear probe reporting
%   balanced accuracy / AUROC for each protected attribute.
% - CONTROLS (name what they control for, not just that they exist): the
%   random-init baseline and the permutation null separate genuine task-learned
%   demographic signal from trivial confounds -- chiefly body extent (the
%   generic encoder's ~90% body-extent confound). Report the probe relative to
%   these baselines.
